\documentclass[10pt]{article}
\usepackage[utf8]{inputenc}
\usepackage[activeacute,spanish]{babel}
\usepackage[left=1.5cm,top=1.5cm,right=1.5cm, bottom=1.5cm,letterpaper, includeheadfoot]{geometry}

\usepackage{amssymb, amsmath, amsthm}
\usepackage{graphicx}
\usepackage{hyperref}
\usepackage{lmodern,url}
\usepackage{paralist} %util para listas compactas
\usepackage{xcolor}
\usepackage{bbm}
\usepackage{mathrsfs}

%========PAQUETES AGREGADOS===========
%Pseudocodigo
\usepackage{pseudocode}
\usepackage[portuguese, boxruled]{algorithm2e}
\usepackage{wrapfig}
\usepackage{multicol}
\usepackage{graphicx}
\usepackage{caption}
\usepackage{subcaption}
%\captionsetup[table]{labelformat=empty}
\captionsetup[subfigure]{labelformat=empty}
\usepackage{cancel}
%====================================

\usepackage{fancyhdr}
\pagestyle{fancy}
\fancypagestyle{plain}{%
\fancyhf{}
\lhead{\footnotesize\itshape\bfseries\rightmark}
\rhead{\footnotesize\itshape\bfseries\leftmark}
}


% macros
\newcommand{\Q}{\mathbb Q}
\newcommand{\R}{\mathbb R}
\newcommand{\N}{\mathbb N}
\newcommand{\Z}{\mathbb Z}
\newcommand{\C}{\mathbb C}
\newcommand{\BigO}{\mathcal{O}}
%Teoremas, Lemas, etc.
\theoremstyle{plain}
\newtheorem{teo}{Teorema}
\newtheorem{lem}{Lema}
\newtheorem{prop}{Proposición}
\newtheorem{cor}{Corolario}
\newtheorem{obs}{Observación}
\newtheorem{ej}{Ejemplo}
\renewcommand{\qedsymbol}{\rule{0.7em}{0.7em}}
\renewenvironment{proof}{{\bfseries \noindent Demostración}}{ \qed \\}


\theoremstyle{definition}
\newtheorem{defi}{Definición}
% fin macros


\newcommand{\catnum}{5} %numero de catedra
\newcommand{\fecha}{13 de Septiembre 2016 }

%%%%%%%%%%%%%%%%%%

%Macros para este documento
\newcommand{\cin}{\operatorname{cint}}



\begin{document}
%Encabezado
\fancyhead[L]{Facultad de Ciencias Físicas y Matemáticas}
\fancyhead[R]{Universidad de Chile}
\vspace*{-1.2 cm}
\begin{minipage}{0.6\textwidth}
\begin{flushleft}
\hspace*{-0.5cm}\textbf{MA3402-1 Estadística. Primavera 2016}\\
\hspace*{-0.5cm}\textbf{Profesor:} Raul Gouet\\
\hspace*{-0.5cm}\textbf{Escriba:} Manuel Cáceres\\
\hspace*{-0.5cm}\textbf{Fecha:} \fecha
\end{flushleft}
\end{minipage}
\begin{minipage}{0.36\textwidth}
\begin{flushright}
\includegraphics[scale=0.3]{imagenes/fcfm_dcc}
\end{flushright}
\end{minipage}
\bigskip
%Fin encabezado

\begin{center}
\LARGE\textbf{Clase \catnum}
\end{center}
\section{Notas Finales sobre EIVUM (Teorema de Rao-Blackwell)}
\begin{enumerate}
\item Se considera $g(\theta)$ estimable insesgadamente
\item Hay un estimador insesgado ``inicial'' $\tilde{g}(X)$ para $g(\theta)$ (de cuadrado integrable)
\item Se tiene un estadístico suficiente $T$
\item El TRB permite mejorar $\tilde{g}$ condicionando en $T$, $\hat{g}=\mathbb{E}_{\theta}(\tilde{g}(X)|T)$
\item Está el concepto de completitud de $T$ (se refiere a la familia de leyes de $\mathbb{P_{\theta, T}}$) ssi no existen estimadores insesgados de 0, no triviales, función de $T$.
\item Si $T$ es suficiente completo, entonces existe una única función de $T$ que es estimador insesgado de $g(\theta)$. La idea es encontrarla, incluso sin usar TRB
\item El estimador $\hat{g}(T)$ en 5. es el EIVUM
\item La función de ``pérdida'' $(\cdot-g(\theta))^2$ en el TRB se puede reemplazar por otra función convexa
\item Si aplicamos TRB cuando $T$ no es completo, entonces resulta algo mejor, pero sera mejorable
\item Sobre completitud y suficiencia minimal. Si $T$ es completo, entonces minimal (Bahadur)
\end{enumerate}
\section{Cota de Cramer-Rao}
Este es un resultado elegante que entrega una cota (inferior) para la varianza de un estimador insesgado de alguna función $g(\theta)$.\\
Aquí se necesitan buenas ``condiciones de regularidad'' en términos de la dependencia de $f_{\theta}(x)$ del parámetro $\theta$.\\
Suponemos que se dispone de densidades para $f_{\theta}$ para $X$(objeto aleatorio).\\
Sea $\tilde{g}(X)$ estimador insesgado de $g(\theta)$. Entonces:
\begin{align*}
\mathbb{E}_{\theta}(\tilde{g}(x)) &= \int_{\mathfrak{X}}\tilde{g}(X)f_{\theta}(x) dx
\end{align*}
Supongamos $g(\theta)$ derivable con respecto a $\theta$ y denotemos por $g'(\theta)$ la derivada. Es decir,\footnote{Suponemos condiciones que garantizan pasar la derivada dentro. En algunos ejemplos clásicos esto no funciona.}
\begin{align*}
g'(\theta) &= \frac{\partial}{\partial \theta} \int_{\mathfrak{X}} \tilde{g}(x)f_{\theta}(x)dx\\
&= \int_{\mathfrak{X}}\tilde{g}(x) \frac{\partial f_{\theta}(x)}{\partial \theta} dx\\
&= \int_{\mathfrak{X}}\left(\tilde{g}(x) \frac{\partial f_{\theta}(x)}{\partial \theta}\frac{1}{f_{\theta}(x)}\right)f_{\theta}(x) dx\\
&= \int_{\mathfrak{X}}\frac{\partial \log f_{\theta}(x)}{\partial \theta} f_{\theta}(x) dx
\end{align*}
Reescribiendo lo anterior tenemos $\mathbb{E}_{\theta}\left(\tilde{g}(X)\frac{\partial \log f_{\theta}(x)}{\partial \theta}(X)\right) = g'(\theta)$.\\

Por otra parte
\begin{align*}
1 &= \int_{\mathfrak{X}}f_{\theta}(x) dx
\end{align*}
Que si la derivamos con respecto a $\theta$ y suponemos condiciones de regularidas, obtenemos que
\begin{align*}
0 &= \int_{\mathfrak{X}} \frac{\partial f_{\theta}(x)}{\partial \theta} dx = \int_{\mathfrak{X}}\frac{\partial \log f_{\theta}(x)}{\partial \theta} f_{\theta}(x) dx
\end{align*}
Por lo que tenemos que $\mathbb{E}\left(\frac{\partial \log f_{\theta}(X)}{\partial \theta}\right) = 0$\\
Para obtener la cota de CR usamos la desigualdad de Cauchy-Shwarz($<x,y>^2 \le ||x||^2||y||^2$).\\
Aquí el producto interno es la integral de un producto de funciones.\\
Notemos que :
\begin{align*}
g'(\theta) &= \mathbb{E}_{\theta}\left((\tilde{g}(X)-g(\theta)) \frac{\partial \log f_{\theta}(x)}{\partial \theta}\right)\\
&= Cov_{\theta}\left(\tilde{g}(X), \frac{\partial \log f_{\theta}(x)}{\partial \theta}\right)
\end{align*}
Cauchy-Shwarz se traduce en que $Cov^2 \le Var_{1} \cdot Var_{2}$, con lo que llegamos a que :
\begin{align*}
&(g'(\theta))^2 \le \mathbb{V}_{\theta}(\tilde{g}(X)) \mathbb{V}_{\theta}\left(\frac{\partial \log f_{\theta}(x)}{\partial \theta}\right)\\
\Rightarrow & \mathbb{V}_{\theta}(\tilde{g}(X)) \ge \frac{(g'(\theta))^2}{I_{\theta}}
\end{align*}
\begin{defi}(Información de Fisher)\\
Información que $X$ ``contiene'' sobre $\theta$ (bajo las condiciones necesarias de regularidad se cumple la última igualdad)
\begin{align*}
I_{\theta} &= \mathbb{V}_{\theta}\left(\frac{\partial \log f_{\theta}(x)}{\partial \theta}\right) \\
&= \mathbb{E}_{\theta}\left(\left(\frac{\partial \log f_{\theta}(x)}{\partial \theta}\right)^2\right)\\
&= -\mathbb{E}_{\theta}\left(\frac{\partial^2 \log f_{\theta}(x)}{\partial \theta^2}\right)
\end{align*}
\end{defi}
¿Se alcanza la cota de CR?\\

En casos muy especiales. Aprovechando est cota podemos definir la ``eficiencia'' de $\tilde{g}(X)$, simplemente comparando la varianza con la cota :
\begin{align*}
0 \le ef_{\theta}(\tilde{g}) = \frac{\frac{(g'(\theta))^2}{I_{\theta}}}{\mathbb{V}_{\theta}(\tilde{g}(X))} \le 1
\end{align*}
Un estimador que alcanza la cota se llama estimador de Cramér-Rao.

\begin{obs}
Si hay dependencia lineal, es decir si $\frac{\partial \log f_{\theta}(x)}{\partial \theta} = a(\theta)(\tilde{g}(X)-g(\theta))$, entonces se alcanza la cota.
\end{obs}

\begin{ej} Modelo Bernoulli\\
\begin{align*}
f_{\theta}(x) &= \theta^{\sum{x_{i}}(1-\theta)^{n-\sum{x_{i}}}}\\
\log f_{\theta}(x) = \sum x_{i} \log{\theta} + (n-\sum{x_{i}})\log{(1-\theta)}\\
\frac{\partial \log f_{\theta}(x)}{\partial \theta} &= \frac{\sum x_{i}}{\theta} - \frac{n-\sum{x_{i}}}{1-\theta}\\
&= \frac{(1-\theta)\sum{x_{i}}-\theta (n-\sum{x_{i}})}{\theta (1-\theta)}\\
&= \frac{n(\overline{X}-\theta)}{\theta (1-\theta)}
\end{align*}
\underline{Conclusión:} $\overline{X}$ es estimador insesgado de CR para $g(\theta)= 0 \theta$ y $\mathbb{V}(\overline{X}) = \frac{\theta (1-\theta)}{n}$ y además $I_{\theta} = \frac{n}{\theta (1-\theta)}$
\end{ej}

\begin{ej} Modelo Poisson de parámetro $\theta$, sobre una MAS $X=(X_{1},X_{2},\ldots,X_{n})$\\
\begin{align*}
\mathbb{P}_{\theta}(X_{i}=k) &= \frac{e^{-\theta}\theta^{k}}{k!} &\ \theta >0,\ k\in \{0,1,\ldots\}\\
f_{\theta}(x) &= \prod_{i=1}^n f_{\theta}(x_{i}) = \frac{e^{-n\theta}\theta^{\sum x_{i}}}{\prod x_{i}!}&\\
\log f_{\theta}(x) &= -n\theta + \sum x_{i} \log \theta - \sum{\log x_{i}!}&\\
\frac{\partial \log f_{\theta}(x)}{\partial \theta} &= -n + \frac{\sum x_{i}}{\theta} = \frac{n}{\theta}(\overline{X}-\theta)\\
I_{\theta} = \frac{n}{\theta} & \overline{X}, \text{estimador insesgado de CR} & g(\theta) = \theta
\end{align*}
Recordamos el caso de $g(\theta) = e^{-\theta}$, donde tenemos un EIVUM $\left(1-\frac{1}{n}\right)^{n\overline{X}}$.\\
La cota para cualquier estimador insesgado de $g(\theta) = e^{-\theta}$ es $\mathbb{V}_{\theta}(\tilde{g}(X)) \ge \frac{e^{-2\theta}}{I_{\theta}} = \frac{\theta e^{-2\theta}}{n}$. Se comprueba que el EIVUM no la alcanza.
\end{ej}

\section{Estimación de máxima verosimilitud}
Es un primer método y más conocido método para encontrar (no necesariamente inesgado).\\
Bienvenido porque hasta aquí solo postulábamos estimadores (salidos de la intuición).\\
La función de verosimilitud no es otra cosa que la densidad $f_{\theta}(x)$, pero vista como función de $\theta$.\\
Se define $L(\theta) = f_{\theta}(x)$.\\
Vamos a maximizar $L$ en $\theta$ y así definimos las soluciones que se llaman estimadores de máxima verosimilitud.\\
Se define el(los) estimador(es) de máxima verosimilitud de $\theta$, como las soluciones al problema $max_{\theta \in \Theta}L(\theta)$ (no tenemos garantía de existencia de máximo).\\
La solución del problema es un conjunto (no vacío), ya sea $arg\max_{\theta \in \Theta} L(\theta)$ o  $arg\sup_{\theta \in \Theta} L(\theta)$.\\
Diremos que $\hat{\theta}(x)$ es un estimador de máxima verosimilitud de $\theta$ si $\hat{\theta}(x) \in arg\max_{\theta \in \Theta} L(\theta) \Leftrightarrow L(\hat{\theta}(X)) \ge L(\theta) \forall \theta \in \Theta$.
\begin{ej} Modelo Bernoulli\\
\begin{align*}
L(\theta) = \theta^{\sum{x_{i}}}(1-\theta)^{n-\sum{x_{i}}}\\
\max_{\theta \in \Theta = \left[0,1\right]} L(\theta)
\end{align*}
,por comodidad operamos con $l(\theta) = \log{L(\theta)}$.

\begin{align*}
l(\theta) &= \sum x_{i} \log{\theta} + (n-\sum x_{i}) \log{(1-\theta)}\\
\frac{\partial l(\theta)}{\partial \theta} &= \frac{n(\overline{X}-\theta}{\theta (1-\theta)} = 0\\
\Rightarrow &\hat{\theta}(X) = \overline{X}
\end{align*}
Que es el único punto estacionario, que es máximo y por lo tanto estimado de máxima verosimilitud (EMV)
\end{ej}
\end{document}